% Generated draft: requires mathematical and source review.
\documentclass[11pt]{article}
\usepackage[margin=25mm]{geometry}
\usepackage{fontspec}
\setmainfont{Latin Modern Roman}
\usepackage{amsmath,amssymb,mathtools,unicode-math}
\setmathfont{Latin Modern Math}
\usepackage{enumitem,needspace,xurl,hyperref}
\setlength{\parindent}{0pt}
\setlength{\parskip}{5pt}
\setlength{\emergencystretch}{3em}
\newcommand{\N}{\mathbb{N}}
\newcommand{\Z}{\mathbb{Z}}
\newcommand{\Q}{\mathbb{Q}}
\newcommand{\R}{\mathbb{R}}
\newcommand{\C}{\mathbb{C}}
\newcommand{\F}{\mathbb{F}}
\newcommand{\A}{\mathbb{A}}
\newcommand{\PP}{\mathbb{P}}
\newcommand{\E}{\mathbb{E}}
\newcommand{\eps}{\varepsilon}
\newcommand{\dd}{\,\mathrm{d}}
\begin{document}
\section*{UnsolvedMath ID: 8800066}
\textbf{Source problem code:} catalogue entry
\subsection{Definitions and mathematical statement}
A rank-$n$ lattice is $L=B\mathbb Z^n$ for a linearly independent basis matrix $B$, with a specified norm and finite input encoding. A cyclic lattice in $\mathbb Z^n$ is invariant under cyclic coordinate shift. An ideal lattice is obtained from an ideal in a specified number-field order or quotient polynomial ring through a specified coefficient or canonical embedding; the ring, dimension, and induced norm are part of its computational definition.

Construct efficient reductions from computational problems on arbitrary lattices to corresponding problems on cyclic or ideal lattices, with controlled parameter losses. For example, $\gamma$-approximate shortest vector asks for nonzero $v\in L$ satisfying
\[
\|v\|\leq\gamma\lambda_1(L),\qquad
\lambda_1(L)=\min_{0\ne w\in L}\|w\|.
\]
An arbitrary-lattice-to-ideal-lattice reduction for this example would use a solver for an explicitly specified structured instance to recover a solution of the original instance, with quantified rank, input-length, and approximation-factor blowup. Merely embedding an abstract lattice without controlling its geometry is not such a reduction.

The source does not select the task, exact versus approximation regime, allowed rings, embedding/norm, or deterministic/randomized/quantum reduction model. Its request is broader than SVP alone. The summary reports worst-case-to-average-case reductions starting from already structured lattices; those do not by themselves reduce arbitrary-lattice worst-case problems to structured-lattice worst-case problems, so they are retained only as related partial progress.
\subsection{Short English statement}
Reduce problems on general lattices to analogous problems on cyclic or ideal lattices, while controlling dimension, input size, norm distortion, and approximation loss. Existing structured-to-average-case reductions address a different reduction direction.
\subsection{Sources}
\begin{itemize}
\item Daniele Micciancio, IPAM \emph{Workshop on Number Theory and Cryptography --- Open Problems} (2006), attachment 6240, PDF page 36, schedule source \url{https://www.ipam.ucla.edu/programs/workshops/workshop-i-number-theory-and-cryptography-open-problems/?tab=schedule}; accessed 2026-07-29. The supplied provenance records a PDF/legacy-PPT visual audit.
\item Peikert--Rosen and Lyubashevsky--Micciancio are the related worst-case/average-case reduction attributions in the supplied summary, which also mentions ring-LWE/PLWE without full publication identifiers or exact hypotheses.
\item Local registry \texttt{data/unsolvedmath/problems.json}, permanent UnsolvedMath ID $8800066$, source code catalogue entry, \emph{Reducing arbitrary lattices to ideal lattices}. Dataset attribution: Ulam AI and the UnsolvedMath Contributors.
\end{itemize}
\end{document}
